\documentclass[12pt]{article}
\usepackage[utf8]{inputenc}
\usepackage[spanish]{babel}
\usepackage{listings}

\usepackage{graphicx}



\title{Informe I (Arquitectura del Computador)}
\author{Jesus Contreras (C.I: 31.021.750)}
\date{2026}

\begin{document}
	
	\maketitle
	
	\section{Implementaciones en MIPS32}
	

	\subsection{¿Cómo se implementa la recursividad en MIPS32? ¿Qué papel cumple la pila (sp)?}
	En MIPS32, la recursividad se implementa utilizando la pila (apuntada por el registro sp) para guardar el contexto de cada llamada. En el código proporcionado, podemos observar:
	
	\begin{itemize}
		\item Al inicio de paridad\_recursiva, se reservan 8 bytes en la pila: subi \$sp,\$sp, 8
		\item Se guardan el registro de retorno \$ra y el argumento \$a0 en la pila: sw \$ra, 4(\$sp) y sw \$a0, 0(\$sp)
		\item Esto permite que después de la llamada recursiva, se puedan restaurar estos valores: lw \$a0, 0(\$sp) y lw \$ra, 4(\$sp)
		\item Finalmente se libera el espacio de pila: addi \$sp, \$sp, 8
	\end{itemize}
	
	La pila es fundamental porque cada llamada recursiva tiene su propio conjunto de variables locales y dirección de retorno. Sin ella, no sería posible mantener la independencia de cada nivel recursivo.
	
	\subsection{¿Qué riesgos de desbordamiento existen? ¿Cómo mitigarlos?}
	El principal riesgo es el stack overflow (desbordamiento de pila), que ocurre cuando:
	\begin{itemize}
		\item La profundidad de la recursión es muy grande (ej. n = 100000)
		\item Cada llamada consume 8 bytes en la pila, por lo que para n=10000 se usarían 80KB aproximadamente
		\item En MARS, la pila tiene un tamaño limitado (típicamente 1MB)
	\end{itemize}
	
	Para mitigarlo:
	\begin{itemize}
		\item Usar enfoques iterativos cuando se prevean números grandes
		\item Verificar que los casos base estén bien definidos para evitar recursión infinita
		\item En sistemas críticos, calcular el máximo de llamadas posibles y dimensionar la pila adecuadamente
		\item Implementar técnicas como recursión de cola (tail recursion) cuando sea posible
	\end{itemize}
	
	\subsection{¿Qué diferencias encontraste entre una implementación iterativa y una recursiva en cuanto al uso de memoria y registros?}
	\begin{itemize}
		\item \textbf{Uso de memoria:} 
		\begin{itemize}
			\item Recursiva: Utiliza la pila dinámicamente. Cada llamada consume 8 bytes adicionales.
			\item Iterativa: No usa memoria dinámica. Solo utiliza registros (\$t0, \$t1, \$t2) y una variable en \$s0.
		\end{itemize}
		\item \textbf{Registros:}
		\begin{itemize}
			\item Recursiva: Requiere guardar y restaurar \$ra y \$a0 constantemente. Hay mayor movimiento de datos entre registros y memoria.
			\item Iterativa: Todo el procesamiento se hace con registros temporales. No hay acceso a memoria más allá de la inicial.
		\end{itemize}
	\end{itemize}
	
	\subsection{¿Qué diferencias encontraste entre los ejemplos académicos del libro y un ejercicio completo y operativo en MIPS32?}
	Los ejemplos académicos suelen presentar:
	\begin{itemize}
		\item Fragmentos de código aislados y perfectamente comentados
		\item Suposiciones de que los datos ya están en registros
		\item Omisión de la interacción con el usuario
	\end{itemize}
	
	Un ejercicio completo como este implica:
	\begin{itemize}
		\item Manejo de entrada/salida con syscalls (imprimir mensajes, leer enteros)
		\item Gestión correcta de la pila en funciones no triviales
		\item Preservación de registros según las convenciones de MIPS
		\item Pruebas con distintos valores y casos límite
		\item Depuración real en simuladores (MARS) donde se detectan errores de desbordamiento, uso incorrecto de registros, etc.
	\end{itemize}
	
	\subsection{Elaborar un tutorial de la ejecución paso a paso en MARS}
	\begin{enumerate}
		\item \textbf{Abrir MARS} y cargar el archivo \texttt{paridad recursiva.asm} (o el iterativo).
		\item \textbf{Ensamblar} el código: Run $\rightarrow$ Assemble (F3).
		\item En la pestaña Execute, observar el código ensamblado y los segmentos de datos.
		\item \textbf{Establecer breakpoints:}
		\begin{itemize}
			\item En la línea jal paridad\_recursiva (para ver la primera llamada)
			\item En la línea subi \$a0, \$a0,1 (para ver cada llamada recursiva)
		\end{itemize}
		\item \textbf{Ejecutar:} Run $\rightarrow$ Go (F5) e ingresar un número pequeño (ej. 3).
		\item Usar Single Step (F7) para avanzar instrucción por instrucción.
		\item \textbf{Observar:}
		\begin{itemize}
			\item Cómo cambia \$sp (stack pointer) en cada llamada
			\item Los valores guardados en la pila (ventana Data Segment)
			\item El registro \$ra antes y después de jal
			\item El resultado final en \$v0
		\end{itemize}
		\item Para la versión iterativa, observar el bucle for y cómo cambian \$t0 y \$t1.
	\end{enumerate}
	
	\subsection{Justificar la elección del enfoque (iterativo o recursivo) según eficiencia y claridad en MIPS}
	\begin{itemize}
		\item \textbf{Enfoque recursivo:}
		\begin{itemize}
			\item Claridad: Muy claro conceptualmente, refleja directamente la definición matemática paridad(n) = 1 - paridad(n-1).
			\item Eficiencia: Baja. Cada llamada implica overhead de pila (8 bytes) y acceso a memoria. Para n=1000, se realizan 1000 llamadas y se usan 8KB de pila.
			\item Riesgo: Desbordamiento para n grandes.
		\end{itemize}
		\item \textbf{Enfoque iterativo:}
		\begin{itemize}
			\item Claridad: También claro, aunque la lógica de alternancia (result = 1 - result) puede no ser intuitiva a primera vista.
			\item Eficiencia: Alta. Usa solo registros, sin acceso a memoria. Realiza exactamente n iteraciones con operaciones simples.
			\item Robustez: No hay riesgo de desbordamiento de pila.
		\end{itemize}
	\end{itemize}
	
	\textbf{Conclusión:} Para este problema específico, el enfoque \textbf{iterativo es superior} en eficiencia y seguridad, manteniendo una claridad aceptable. El recursivo se justifica principalmente con fines didácticos para aprender manejo de pila y recursividad en ensamblador.
	
	\subsection{Análisis y Discusión de los Resultados}
	\begin{itemize}
		\item Ambos algoritmos producen el mismo resultado para cualquier entrada positiva. El recursivo implementa correctamente la fórmula, mientras que el iterativo simula el mismo comportamiento mediante un bucle.
		\item \textbf{Rendimiento:} Al probar con números grandes (ej. n=10000), el recursivo puede generar desbordamiento de pila en MARS, mientras que el iterativo lo maneja sin problemas.
		\item \textbf{Depuración:} La versión recursiva es más difícil de depurar debido al gran número de llamadas y la gestión de pila. La versión iterativa permite seguir fácilmente el flujo del programa.
		\item \textbf{Uso de registros:} El recursivo usa \$a0, \$ra y \$sp intensivamente; el iterativo usa \$t0, \$t1, \$t2 y \$s0. El iterativo respeta mejor la convención de uso de registros en MIPS.
	\end{itemize}
	
\end{document}
